\documentclass[10pt,a4paper]{amsart}
\usepackage[margin=19mm]{geometry}
\usepackage{amsmath,amssymb,mathtools}
\usepackage[scr=rsfs,cal=euler]{mathalfa}
\usepackage{xfrac}
\usepackage[unicode,hidelinks,bookmarks=false]{hyperref}
\newcommand{\ie}{i.e.\ }
\newcommand{\cf}{cf.\ }
\newcommand{\resp}{respectively\ }
\newcommand{\Subsection}{\subsection}
\providecommand{\nobreakdash}{\nobreak\hbox{-}}
\setlength{\emergencystretch}{2em}
\multlinegap0pt
\usepackage{enumerate}
\newtheorem{lemma}{Lemma}
\newtheorem{theorem}{Theorem}
\title{The curvature of convex sum of metrics and applications}
\author{Leonardo F. Cavenaghi, Giovane Galindo, Llohann D. Speran\c{c}a}
\date{Source: arXiv:2106.14781v2 (CC BY 4.0)}
\begin{document}
\maketitle

\noindent\emph{Source and licence.} The curvature of convex sum of metrics and applications. Version arXiv:2106.14781v2 (CC BY 4.0), distributed by arXiv under the Creative Commons Attribution 4.0 International licence, http://creativecommons.org/licenses/by/4.0/, as recorded in the arXiv licence metadata for this exact version and shown on the abstract page https://arxiv.org/abs/2106.14781v2. Full licence notice: \texttt{licenses/arxiv-2106.14781-CC-BY-4.0.txt}.\par\smallskip

\noindent\textit{Editorial scope.} Counted unit: the article's theorem thm:rvariation (source0.tex lines 324--340). The article's complete original proof of it is delivered with it (source0.tex lines 342--372, one \texttt{proof} environment of 31 lines). No mathematics is written, repaired or re-derived: the statement and the proof are the article's own lines, in the article's own notation, and no retained line is substituted or re-flowed.  Retained uncounted as the source-local context the proof needs: lines 271--283.  Nothing was omitted from any retained span, so the package carries no omission record.  No retained line contains a citation, so the wrapper carries no bibliography and no bibliography entry.  No retained line contains a figure environment, an image-inclusion command or drawing code, so no image is delivered and the package declares no asset dependency.  Editorial formatting only, in addition: an AMS \texttt{amsart} preamble that restates the article's own declarations and macros used by the retained lines, namely the environments \texttt{lemma}, \texttt{theorem} on separate counters, together with the standard packages those lines need.  The retained environments print in this document as Lemma 1, Theorem 1 in order of appearance, whereas the article prints the same items with the numbering of its own full document.  The complete native source of the article as distributed in the arXiv e-print for this exact version is bundled unchanged as \texttt{sources/113157/source0.tex}.
\begin{lemma}\label{lem:resolvetudo}
Assume that $\nabla^0_X X = \nabla^0_X Y = \nabla^0_Y Y = 0$ on $\mathbb{T}^2$. Denote by $\nabla^t$ the Levi-Civita connection of $g_t$. Then for sufficiently small $t$ (such that $\|t(I-P)\| < 1$), the Riemann curvature component admits the expansion
\begin{align}
R_t(X,Y,Y,X) =& (1-t)R_0(X,Y,Y,X) + tR_1(X,Y,Y,X) \nonumber \\
&+ t\Bigl[g_1\bigl((\nabla^1_X X)^\perp, (\nabla^1_Y Y)^\perp\bigr) 
     - g_1\bigl((\nabla^1_X Y)^\perp, (\nabla^1_X Y)^\perp\bigr)\Bigr] \nonumber \\
&+ t\Bigl[g_1\bigl((\nabla^1_X X)^\top, (\nabla^1_Y Y)^\top\bigr) 
     - g_1\bigl((\nabla^1_X Y)^\top, (\nabla^1_X Y)^\top\bigr)\Bigr] \nonumber \\
&- t^2\sum_{n=1}^\infty \Bigl[g_1\bigl(P(t(I-P))^{n-1}\nabla^1_X X, \nabla^1_Y Y\bigr) \nonumber \\
&\qquad\qquad - g_1\bigl(P(t(I-P))^{n-1}\nabla^1_X Y, \nabla^1_X Y\bigr)\Bigr],
\end{align}
where $^\top$ and $^\perp$ denote the $g_1$-orthogonal projections onto $T\mathbb{T}^2$ and its orthogonal complement, respectively, and $P$ is defined by $g_0(P\cdot,\cdot)=g_1(\cdot,\cdot)$.
\end{lemma}

\begin{theorem}\label{thm:rvariation}
Assume that $\nabla^0_X X = \nabla^0_X Y = \nabla^0_Y Y = 0$ on $\mathbb{T}^2$. Denote $R'_0(X,Y,Y,X) := \left.\frac{d}{dt}\right|_{t=0} R_t(X,Y,Y,X)$. Then
\begin{multline}\label{eq:first-variation}
\int_{\mathbb{T}^2} \frac{1}{\|X\wedge Y\|_1^2} R'_0(X,Y,Y,X) \, dA \\
= \int_{\mathbb{T}^2} \frac{1}{\|X\wedge Y\|_1^2} \Bigl[
g_1\bigl((\nabla^1_X X)^\top, (\nabla^1_Y Y)^\top\bigr) 
- g_1\bigl((\nabla^1_X Y)^\top, (\nabla^1_X Y)^\top\bigr) \Bigr] dA,
\end{multline}
where $\|X\wedge Y\|_1^2 := g_1(X,X)g_1(Y,Y) - g_1(X,Y)^2$.

Moreover, if $R'_0(X,Y,Y,X) = 0$, then for any $r \geq 2$,
\begin{equation}\label{eq:higher-order-condition}
\int_{\mathbb{T}^2} \left.\frac{d^r}{dt^r}\right|_{t=0} R_t(X,Y,Y,X) \, dA > 0 
\quad\Longleftrightarrow\quad 
\int_{\mathbb{T}^2} S_r^P(X,Y) \, dA < 0.
\end{equation}
\end{theorem}

\begin{proof}
Differentiating the expansion in Lemma~\ref{lem:resolvetudo} at $t=0$ yields:
\begin{align*}
R'_0(X,Y,Y,X) &= R_1(X,Y,Y,X) - R_0(X,Y,Y,X) \\
&\quad + g_1\bigl((\nabla^1_X X)^\perp, (\nabla^1_Y Y)^\perp\bigr) 
   - g_1\bigl((\nabla^1_X Y)^\perp, (\nabla^1_X Y)^\perp\bigr) \\
&\quad + g_1\bigl((\nabla^1_X X)^\top, (\nabla^1_Y Y)^\top\bigr) 
   - g_1\bigl((\nabla^1_X Y)^\top, (\nabla^1_X Y)^\top\bigr).
\end{align*}
Since $\mathbb{T}^2$ is flat with respect to $g_0$, we have $R_0(X,Y,Y,X)=0$. Moreover, by the Gauss equation for the submanifold $\mathbb{T}^2 \subset M$ with the metric $g_1$, we have:
\[
R_1(X,Y,Y,X) = K^1_{\mathbb{T}^2} \|X\wedge Y\|_1^2 + \left[g_1((\nabla^1_X X)^\perp, (\nabla^1_Y Y)^\perp) - g_1((\nabla^1_X Y)^\perp, (\nabla^1_X Y)^\perp)\right],
\]
where $K^1_{\mathbb{T}^2}$ is the intrinsic Gaussian curvature of $\mathbb{T}^2$ with respect to $g_1$. Substituting this gives:
\[
R'_0(X,Y,Y,X) = K^1_{\mathbb{T}^2} \|X\wedge Y\|_1^2 + g_1((\nabla^1_X X)^\top, (\nabla^1_Y Y)^\top) - g_1((\nabla^1_X Y)^\top, (\nabla^1_X Y)^\top).
\]
Integrating over $\mathbb{T}^2$ and applying the Gauss-Bonnet theorem ($\int_{\mathbb{T}^2} K^1_{\mathbb{T}^2} \, dA = 0$) yields \eqref{eq:first-variation}.

Now assume $R'_0(X,Y,Y,X)=0$. From Lemma~\ref{lem:resolvetudo}, for $r \geq 2$, the $r$th derivative at $t=0$ of the sum $ -t^2\sum_{n=1}^\infty \left[g_1(P(t(I-P))^{n-1}\nabla^1_X X, \nabla^1_Y Y) - g_1(P(t(I-P))^{n-1}\nabla^1_X Y, \nabla^1_X Y)\right] $ receives contributions only from the term where the power of $t$ is exactly $r$. This occurs when $n+1 = r$, i.e., $n = r-1$. For this term,
\begin{align*}
&\left.\frac{d^r}{dt^r}\right|_{t=0} \left( -t^2 \cdot \left[g_1(P(t(I-P))^{r-2}\nabla^1_X X, \nabla^1_Y Y) - \cdots \right] \right) \\
&= \left.\frac{d^r}{dt^r}\right|_{t=0} \left( -t^{r} \left[g_1(P(I-P)^{r-2}\nabla^1_X X, \nabla^1_Y Y) - \cdots \right] \right) \\
&= -r! \left[g_1(P(I-P)^{r-2}\nabla^1_X X, \nabla^1_Y Y) - g_1(P(I-P)^{r-2}\nabla^1_X Y, \nabla^1_X Y)\right].
\end{align*}
This simplifies to:
\[
\left.\frac{d^r}{dt^r}\right|_{t=0} R_t(X,Y,Y,X) = -r! \cdot S_r^P(X,Y).
\]
Integrating over $\mathbb{T}^2$ yields the equivalence \eqref{eq:higher-order-condition}.
\end{proof}

\end{document}
